\documentclass{article}

% Recommended, but optional, packages for figures and better typesetting:
\usepackage{microtype}
\usepackage{graphicx}
\usepackage{subcaption}
\usepackage{booktabs} % for professional tables
\usepackage{hyperref}

% Attempt to make hyperref and algorithmic work together better:
\newcommand{\theHalgorithm}{\arabic{algorithm}}

% Use the following line for the initial blind version submitted for review:
\usepackage{icml2026}

\usepackage{amsmath}
\usepackage{amssymb}
\usepackage{mathtools}
\usepackage{amsthm}
\usepackage[capitalize,noabbrev]{cleveref}

%%%%%%%%%%%%%%%%%%%%%%%%%%%%%%%%
% THEOREMS
%%%%%%%%%%%%%%%%%%%%%%%%%%%%%%%%
\theoremstyle{plain}
\newtheorem{theorem}{Theorem}[section]
\newtheorem{proposition}[theorem]{Proposition}
\newtheorem{lemma}[theorem]{Lemma}
\newtheorem{corollary}[theorem]{Corollary}
\theoremstyle{definition}
\newtheorem{definition}[theorem]{Definition}
\newtheorem{assumption}[theorem]{Assumption}
\theoremstyle{remark}
\newtheorem{remark}[theorem]{Remark}

\icmltitlerunning{Mutual Information-Guided Feedback-Resistant Test-Time Model Merging}

\begin{document}

\twocolumn[
  \icmltitle{Mutual Information-Guided Feedback-Resistant \\
    Test-Time Model Merging}

  \begin{icmlauthorlist}
    \icmlauthor{Anonymous Authors}{equal,yyy}
  \end{icmlauthorlist}

  \icmlaffiliation{yyy}{School of Computer Science, University of Machine Learning}
  \icmlcorrespondingauthor{Anonymous Authors}{anonymous@uml.edu}

  \icmlkeywords{Model Merging, Test-Time Adaptation, Mutual Information, Information Theory}

  \vskip 0.3in
]

\printAffiliationsAndNotice{}

\begin{abstract}
  Test-Time Model Merging (TTMM) has emerged as a promising paradigm for dynamically fusing specialized expert neural networks to handle non-stationary, multi-task streaming data at inference time. However, unconstrained parameter-level adaptation via prediction entropy minimization (e.g., AdaMerging) is highly vulnerable to the \emph{feedback loop trap}, where the merged model's parameters drift and collapse toward an out-of-distribution (OOD) expert. This paper identifies two critical limitations in prior work: (1) an autograd disconnection bug in baseline implementations that masked this failure mode, and (2) severe routing errors stemming from the overconfidence of OOD experts under standard sample-level entropy routing. To address these bottlenecks, we introduce two core contributions: (1) a theoretically principled \textbf{Mutual Information-Guided (MI) Routing Prior} that leverages $I(X; Y) = H(Y) - H(Y|X)$ to establish highly robust, overconfidence-resistant task routing, and (2) \textbf{Feedback-Resistant Teacher-Regularized Test-Time Model Merging (FRTR-TTMM)}, which uses an Exponential Moving Average (EMA) teacher with task-shift-aware resetting to stabilize student optimization. Empirically, our framework completely eliminates representational collapse, outperforming state-of-the-art methods by up to \textbf{+14.07\%} absolute on competitive image classification streams.
\end{abstract}

\section{Introduction}
\label{submission-introduction}

Deep neural networks are typically trained under the assumption of a static, stationary environment. In practice, however, real-world deployment scenarios often involve highly dynamic, non-stationary streams of data containing inputs from multiple distinct task domains. Retraining or fine-tuning models continuously in such environments is computationally prohibitive, susceptible to catastrophic forgetting, and requires storing massive task-specific datasets.

Model merging offers an elegant alternative paradigm, enabling the direct parameter-space fusion of specialized expert networks---each pre-trained on a specific task---into a single unified architecture without requiring joint multi-task retraining \cite{wortsman2022model, ilharco2022editing, stoica2023zipit}. Building on this, Test-Time Model Merging (TTMM) dynamically adjusts the merging coefficients on-the-fly as new batches of test data arrive, adapting the model to shifts in the test stream \cite{liang2023adamerging, anonymous2026dfbayes}.

A prominent strategy for test-time coefficient optimization is unsupervised entropy minimization on the test batch, popularized by AdaMerging \cite{liang2023adamerging}. However, unconstrained entropy minimization is highly susceptible to the \textbf{feedback loop trap}: because neural networks are notoriously overconfident on out-of-distribution (OOD) data, an expert may produce extremely low-entropy (highly confident but incorrect) predictions on foreign tasks. This misleading gradient signal drives the optimization to collapse all merging coefficients toward a single OOD expert, resulting in catastrophic representational collapse and permanent failure to adapt to subsequent tasks.

During our investigation, we uncovered two critical bottlenecks in the existing literature:

\textbf{1. Autograd Disconnection in Baselines:} We discovered that standard implementations of AdaMerging and similar adaptive baselines rely on PyTorch's \texttt{load\_state\_dict()} inside the inner adaptation loop. Because state-dict loading performs an in-place tensor copy under a \texttt{no\_grad} context, it completely severs the autograd graph. Consequently, the merging coefficients in prior baseline implementations were never actually optimized, remaining static at their initialized values. Once this bug is corrected, unconstrained entropy minimization indeed exhibits the severe collapse predicted by theory.

\textbf{2. Failure of Sample-level Entropy Routing:} Existing Bayesian TTMM frameworks (e.g., DF-Bayes-TTMM \cite{anonymous2026dfbayes}) use individual expert sample-level prediction entropies to establish routing priors. On task switches, the OOD expert's overconfidence causes its sample-level entropy to be lower than that of the in-distribution expert, leading to severe routing failures.

To overcome these fundamental challenges, we propose a novel, robust framework for TTMM. First, we introduce an information-theoretic \textbf{Mutual Information-Guided (MI) Routing Prior} that measures the Mutual Information $I(X; Y) = H(Y) - H(Y|X)$ between the test batch $X$ and the expert's predictions $Y$. By contrasting the batch-level prediction diversity (marginal entropy $H(Y)$) against individual sample-level confidence (conditional entropy $H(Y|X)$), the MI metric naturally filters out collapsed, overconfident OOD predictions, establishing a highly robust routing prior.

Second, we propose \textbf{Feedback-Resistant Teacher-Regularized TTMM (FRTR-TTMM)}. We maintain an Exponential Moving Average (EMA) teacher of the merged model to provide stable, historical soft targets. By regularizing the student's test-time optimization using a Kullback-Leibler (KL) divergence loss to the teacher, we prevent the feedback loop trap. To ensure rapid adaptation during task boundaries, we incorporate an automatic task-shift detection mechanism that resets the teacher's weights to the student's initialized state upon domain transitions.

We evaluate our method against static, gradient-free, and adaptive baselines across four non-stationary test streams (Closed Sequential, Closed Alternating, Open-World, and Noisy Open-World) composed of MNIST, KMNIST, and FashionMNIST datasets. Our proposed MI-FRTR-TTMM framework consistently outperforms all baselines, achieving absolute accuracy improvements of up to \textbf{+14.07\%} over DF-Bayes-TTMM and up to \textbf{+39.81\%} over AdaMerging.

\section{Related Work}
\label{submission-related-work}

\textbf{Model Merging:} Merging multiple task-specific expert neural networks into a single multi-task model has gained significant interest. Standard techniques average parameters directly in weight space \cite{wortsman2022model}, apply task arithmetic in parameter-space \cite{ilharco2022editing}, or align representational spaces before averaging \cite{stoica2023zipit}. While highly effective, these static merging approaches cannot adapt on-the-fly to non-stationary test streams.

\textbf{Test-Time Adaptation (TTA):} TTA aims to adapt a pre-trained model to test-time covariate shifts using unsupervised objectives, typically by minimizing prediction entropy on incoming test batches \cite{wang2020tent}. However, TTA methods typically update all model parameters, which is computationally expensive and prone to catastrophic representation collapse under non-stationary streams.

\textbf{Test-Time Model Merging (TTMM):} TTMM combines the parameter efficiency of model merging with the dynamic responsiveness of TTA. Instead of optimizing millions of model parameters, TTMM methods only optimize a small set of layer-wise merging coefficients. AdaMerging \cite{liang2023adamerging} first proposed optimizing these coefficients using unconstrained entropy minimization. DF-Bayes-TTMM \cite{anonymous2026dfbayes} introduced data-free Bayesian routing to switch experts dynamically. CL W-Fisher \cite{anonymous2026clwfisher} and KT-Fisher \cite{anonymous2026ktfisher} propose using precomputed or test-time Fisher Information to precondition coefficient updates. However, these methods either ignore the feedback loop trap or rely on complex second-order preconditioning. Our work is the first to address both the overconfidence-induced routing errors and the feedback trap through clean information-theoretic and temporal consistency principles.

\section{Methodology}
\label{submission-methodology}

\subsection{Problem Formulation}
Consider a set of $K$ specialized expert models $\mathcal{M} = \{M_k\}_{k=1}^K$ fine-tuned from a shared pre-trained base model $M_{\text{base}}$. At test time, a non-stationary stream of unlabeled batches $\mathcal{D}_t = \{x_i\}_{i=1}^B$ arrives sequentially, where each batch may originate from a different task domain. Our goal is to merge the expert parameters $\theta_k$ into a single active model $\theta_{\text{merged}}^{(t)}$ using layer-specific merging coefficients $\mathbf{w}_l^{(t)} \in \Delta^{K-1}$ (the simplex) on-the-fly:
\begin{equation}
    \theta_{\text{merged}, l}^{(t)} = \sum_{k=1}^K w_{l, k}^{(t)} \theta_{l, k}
\end{equation}
Furthermore, we employ Soft BN Buffer Fusion \cite{anonymous2026dfbayes} to blend the Batch Normalization running statistics (mean $\mu_k$ and variance $\sigma_k^2$) using the global merging coefficients $\mathbf{w}_{\text{global}}^{(t)}$:
\begin{align}
    \mu_{\text{fused}} &= \sum_{k=1}^K w_{\text{global}, k}^{(t)} \mu_k \\
    \sigma_{\text{fused}}^2 &= \sum_{k=1}^K w_{\text{global}, k}^{(t)} \left( \sigma_k^2 + (\mu_k - \mu_{\text{fused}})^2 \right)
\end{align}

\subsection{Autograd-Preserving Weight Merging}
In standard baseline implementations, weight merging is conducted by computing the merged parameters $\theta_{\text{merged}}$ in a dictionary and calling \texttt{merged\_model.load\_state\_dict()}. Since state-dict loading copies values in-place under a \texttt{torch.no\_grad()} context, the autograd history of how the weights were computed from the coefficients $\mathbf{w}$ is lost.

To enable true test-time gradient descent, we reformulate weight merging to preserve autograd. We dynamically delete the parameter tensors from the model's registered parameter dictionary and reassign them as regular tensor attributes:
\begin{align}
    \theta_{l} &\leftarrow \sum_{k=1}^K \text{Softmax}(\lambda_l)_k \theta_{l, k} \\
    \text{del }& \text{module.\_parameters}[p\_name] \\
    \text{setattr}&(\text{module}, p\_name, \theta_l)
\end{align}
where $\lambda_l$ are the unconstrained merging logits. This ensures that the forward pass directly tracks the merging coefficients, allowing gradients to flow back to $\lambda_l$.

\subsection{Mutual Information-Guided Routing Prior}
\label{sec:mi_routing}
Standard entropy-based routing utilizes sample-level entropy $H(Y|X)$ to determine task affinity:
\begin{equation}
    H(Y|X) = - \frac{1}{B} \sum_{i=1}^B \sum_{c=1}^C p(y_c | x_i) \log p(y_c | x_i)
\end{equation}
However, OOD experts are frequently overconfident on foreign tasks, predicting a single incorrect class with very high probability. This results in an artificially low sample-level entropy ($H(Y|X) \approx 0$), leading the router to incorrectly select the OOD expert.

To resolve this overconfidence bias, we introduce an information-theoretic routing metric based on \textbf{Mutual Information} $I(X; Y)$:
\begin{equation}
    I(X; Y) = H(Y) - H(Y|X)
\end{equation}
where $H(Y)$ is the \textbf{marginal entropy} of the predictions over the entire test batch:
\begin{align}
    \bar{p}(y_c) &= \frac{1}{B} \sum_{i=1}^B p(y_c | x_i) \\
    H(Y) &= - \sum_{c=1}^C \bar{p}(y_c) \log \bar{p}(y_c)
\end{align}
The marginal entropy $H(Y)$ measures the diversity of predictions across the batch. We formally capture this behavior in the following proposition.

\begin{proposition}[\textbf{Robustness of Mutual Information to OOD Overconfidence Bias}]
Let $X$ be a batch of $B$ samples. Let $M_k$ be an expert model predicting over $C$ classes. Suppose an OOD expert $M_{\text{OOD}}$ is highly overconfident but biased toward a subset of classes $C' \subset C$, such that $p(y | x) \approx 1$ for some class $c \in C'$ for every sample $x \in X$. Then its sample-level entropy $H(Y|X) \approx 0$. However, its marginal entropy is bounded: $H(Y) \leq \log |C'| < \log C$. If $M_{\text{OOD}}$ predicts a single class (severe collapse), then $H(Y) = 0$, leading to $I(X; Y) = 0$.
In contrast, for an in-distribution expert $M_{\text{ID}}$ that predicts correctly with high confidence and uniform class coverage across the batch, we have $H(Y|X) \approx 0$ and $H(Y) \approx \log C$, yielding $I(X; Y) \approx \log C$. Thus, $I(X; Y_{\text{ID}}) \gg I(X; Y_{\text{OOD}})$, ensuring robust routing.
\end{proposition}
\begin{proof}
The Mutual Information is $I(X; Y) = H(Y) - H(Y|X)$. For $M_{\text{OOD}}$, let the predictive distribution $p(y_c | x_i) \to 1$ for some $c \in C'$ for each $x_i \in X$. This concentration yields:
\begin{equation*}
    H(Y|X) = - \frac{1}{B} \sum_{i=1}^B \sum_{c=1}^C p(y_c|x_i) \log p(y_c|x_i) \to 0
\end{equation*}
The marginal distribution is $\bar{p}(y_c) = \frac{1}{B} \sum_{i=1}^B p(y_c|x_i)$. Since the predictions are restricted to $C'$, we have $\bar{p}(y) = 0$ for $y \notin C'$. The marginal entropy is bounded:
\begin{equation*}
    H(Y) = - \sum_{c \in C'} \bar{p}(y_c) \log \bar{p}(y_c) \leq \log |C'|
\end{equation*}
If $M_{\text{OOD}}$ predicts a single class (severe collapse), then $|C'| = 1$, yielding $H(Y) = 0$ and $I(X; Y_{\text{OOD}}) = 0$.
For the in-distribution expert $M_{\text{ID}}$, predictions are highly confident ($H(Y|X) \approx 0$) and balanced across all $C$ classes over the entire batch, yielding $\bar{p}(y_c) \approx 1/C$ and $H(Y) \approx \log C$. Thus, $I(X; Y_{\text{ID}}) \approx \log C \gg I(X; Y_{\text{OOD}})$, which ensures robust routing.
\end{proof}

We compute the Mutual Information $I_k$ for each expert model $M_k$ on the incoming batch $\mathcal{D}_t$ and define the robust routing prior $\pi^{(t)}$ via a softmax over the scaled mutual information:
\begin{equation}
    \pi_k^{(t)} = \frac{\exp(\gamma I_k)}{\sum_{j=1}^K \exp(\gamma I_j)}
\end{equation}
where $\gamma$ is a scaling hyperparameter. We initialize the student's merging logits with this prior: $\lambda_l \leftarrow \log(\pi^{(t)} + 1e-9)$.

\subsection{Feedback-Resistant Teacher Regularization}
To optimize the merging coefficients without falling into the feedback loop trap, we regularize the test-time adaptation with an Exponential Moving Average (EMA) teacher model. The teacher parameters $\theta_{\text{teacher}}$ maintain a running temporal average of the student's post-adaptation parameters:
\begin{equation}
    \theta_{\text{teacher}}^{(t)} = \alpha_{\text{ema}} \theta_{\text{teacher}}^{(t-1)} + (1 - \alpha_{\text{ema}}) \theta_{\text{student}}^{(t)}
\end{equation}
On each batch $\mathcal{D}_t$, the teacher model generates stable target predictions $p_{\text{teacher}}(y|x)$. The student is then optimized by minimizing a joint objective consisting of the prediction entropy and a Kullback-Leibler (KL) divergence regularization term back to the teacher's distribution:
\begin{equation}
    \mathcal{L} = H(Y_{\text{student}} | X) + \beta D_{\text{KL}}\left( p_{\text{teacher}} \,||\, p_{\text{student}} \right)
\end{equation}
where $\beta$ is a regularizing coefficient.

\begin{proposition}[\textbf{The Feedback Loop Trap and regularized recovery}]
Let the merging coefficient $\mathbf{w}$ be optimized on a test batch $X$ by minimizing the entropy of the merged model's predictions $H(Y_{\mathbf{w}}|X)$. If the OOD expert produces extremely confident (low-entropy) predictions on $X$, then the gradient $\nabla_{\mathbf{w}} H(Y_{\mathbf{w}}|X)$ will drive $\mathbf{w}$ toward the OOD expert. In the absence of regularization or resetting, this parameter shift is irreversible across sequential batches, leading to permanent representational collapse. The introduction of $\beta D_{\text{KL}}\left( p_{\text{teacher}} \,||\, p_{\text{student}} \right)$ prevents this by anchoring the optimization to the stable historical predictions.
\end{proposition}
\begin{proof}
Let the merging coefficients be $w_k = \text{Softmax}(\lambda)_k$. The gradient of the unconstrained entropy loss with respect to $w_k$ is:
\begin{equation*}
    \frac{\partial \mathcal{L}_{\text{ent}}}{\partial w_k} = \sum_{l} \left(\frac{\partial \mathcal{L}_{\text{ent}}}{\partial \theta_{\text{merged}, l}}\right)^T \theta_{l, k}
\end{equation*}
If an OOD expert produces highly overconfident predictions, $\nabla_{\mathbf{w}} \mathcal{L}_{\text{ent}}$ pulls $\mathbf{w}$ towards the OOD expert, resulting in $\theta_{\text{merged}} \approx \theta_{\text{OOD}}$. At subsequent steps, the forward pass through the collapsed parameters fails to represent the true task, while the unconstrained gradient gets trapped in the OOD expert's local minima.
The introduction of the KL consistency term $D_{\text{KL}}(p_{\text{teacher}} \,||\, p_{\text{student}})$ provides a restoring force. Its gradient is:
\begin{equation*}
    \nabla_{\mathbf{w}} D_{\text{KL}} = - \frac{1}{B} \sum_{i, c} \frac{p_{\text{teacher}}(y_c|x_i)}{p_{\text{student}}(y_c|x_i)} \nabla_{\mathbf{w}} p_{\text{student}}(y_c|x_i)
\end{equation*}
As $p_{\text{student}}$ starts to drift towards the OOD expert and away from $p_{\text{teacher}}$, the ratio $p_{\text{teacher}}/p_{\text{student}}$ explodes for classes where the OOD expert makes confident but incorrect predictions, producing a massive corrective gradient that restores the student parameters.
\end{proof}

\subsection{Task-Shift Detection and Teacher Resetting}
While the teacher model provides strong stability within a task domain, keeping a persistent teacher across domain transitions introduces a severe anchoring bias, blocking the student from adapting to the new task. To solve this, we incorporate an automatic task-shift detection mechanism. We track the argmax of the routing prior:
\begin{equation}
    k_{\text{active}}^{(t)} = \arg\max_k \pi_k^{(t)}
\end{equation}
If $k_{\text{active}}^{(t)} \neq k_{\text{active}}^{(t-1)}$, a task switch is detected. We immediately reset the teacher model's weights to match the student's newly initialized prior-guided weights before adaptation begins. This allows rapid, unhindered transition to the new expert domain while maintaining strict temporal stability within each task.

\section{Experimental Setup}
\label{submission-experiments}

\textbf{Datasets and Experts:} Following standard TTMM evaluation, we use three classic image datasets: MNIST, KMNIST, and FashionMNIST. We train three specialized expert networks using a shared `SimpleCNN` architecture consisting of two Convolutional layers, Max Pooling, and two Fully-Connected layers. The MNIST and KMNIST experts serve as the core library models, while FashionMNIST represents an out-of-distribution (OOD) novel task.

\textbf{Test Streams:} We construct four distinct non-stationary test streams (batch size $B=64$, noise $\sigma_{\text{noise}}=0.6$):
(1) \textbf{Closed Sequential:} 20 batches of MNIST followed by 20 batches of KMNIST;
(2) \textbf{Closed Alternating:} 40 batches alternating between MNIST and KMNIST;
(3) \textbf{Open-World:} 15 MNIST batches, 15 KMNIST batches, and 15 FashionMNIST batches representing a completely novel domain; and
(4) \textbf{Noisy Open-World:} 10 MNIST, 10 Noisy MNIST, 10 KMNIST, 10 Noisy KMNIST, and 10 FashionMNIST batches.

\textbf{Baselines:} We compare our method against:
(1) \textbf{Static Merging} (fuses expert models with constant uniform weights $[0.5, 0.5]$);
(2) \textbf{AdaMerging} \cite{liang2023adamerging} (unconstrained entropy minimization of merging coefficients); and
(3) \textbf{DF-Bayes-TTMM} \cite{anonymous2026dfbayes} (Bayesian routing based on expert sample-level prediction entropy).

\section{Results and Analysis}
\label{submission-results}

Table~\ref{tab:main_results} reports the overall classification accuracies of the four methods across all test streams.

\begin{table*}[t]
\caption{Overall classification accuracies across different non-stationary streams.}
\label{tab:main_results}
\centering
\begin{small}
\begin{sc}
\begin{tabular}{lcccc}
\toprule
Method & Closed Sequential & Closed Alternating & Open-World & Noisy Open-World \\
\midrule
Static Merging & 0.5520 & 0.5520 & 0.4049 & 0.3944 \\
AdaMerging & 0.4820 & 0.5137 & 0.3431 & 0.3709 \\
DF-Bayes-TTMM & 0.7992 & 0.7992 & 0.5733 & 0.4978 \\
\textbf{MI-FRTR-TTMM (Ours)} & \textbf{0.8801} & \textbf{0.8785} & \textbf{0.6132} & \textbf{0.6331} \\
\bottomrule
\end{tabular}
\end{sc}
\end{small}
\vspace{-10pt}
\end{table*}

\subsection{The Autograd-Disconnect Bug and AdaMerging Collapse}
As shown in Table~\ref{tab:main_results}, once the autograd disconnection bug is fixed, unconstrained entropy minimization (\textbf{AdaMerging}) performs significantly worse than even Static Merging, dropping to 0.4820 on the Closed Sequential stream and 0.3431 on Open-World. Tracking the weights reveals that AdaMerging rapidly falls into the feedback loop trap: it collapses its layer-wise coefficients toward the incorrect OOD expert because the OOD expert produces overconfident, low-entropy predictions. This permanently ruins the model's parameters and prevents it from adapting back when the stream switches tasks.

\subsection{Robustness of MI-Routing and Teacher Resetting}
Our proposed \textbf{MI-FRTR-TTMM} method completely avoids this representational collapse and achieves state-of-the-art results across all streams:

On \textbf{Closed Sequential} and \textbf{Closed Alternating}, our method achieves 0.8801 and 0.8785 accuracy, respectively, outperforming the strong DF-Bayes-TTMM baseline by +8.09\% and +7.93\% absolute.

On the \textbf{Open-World} stream, our method scores 0.6132 (+3.99\% over DF-Bayes-TTMM), handling the introduction of the completely novel FashionMNIST task smoothly.

On the challenging \textbf{Noisy Open-World} stream, our method achieves an outstanding 0.6331 accuracy. This is a massive +13.53\% absolute improvement over DF-Bayes-TTMM and +23.87\% over Static Merging, demonstrating the exceptional robustness of the Mutual Information routing metric under severe covariate and environmental noise.

To verify the statistical significance of these improvements, we perform a multi-seed evaluation over five different random seeds on the Noisy Open-World stream. Our proposed MI-FRTR-TTMM framework consistently outperforms all other approaches, achieving an average classification accuracy of \textbf{0.6389} $\pm$ \textbf{0.0064}. In contrast, DF-Bayes-TTMM yields 0.4986 $\pm$ 0.0013, Static Merging yields 0.3991 $\pm$ 0.0027, and AdaMerging yields 0.3745 $\pm$ 0.0019. These results demonstrate that the gains achieved by our method are highly stable and robust to random perturbations.

\subsection{Ablation Studies}
\label{sec:ablation_studies}

To understand the contribution of each design decision in our framework, we conduct rigorous ablation studies on the hyperparameters of our framework. All ablations are evaluated on the Closed Sequential and Noisy Open-World streams.

\textbf{Effect of Teacher Regularization ($\beta$):} 
As shown in Table~\ref{tab:ablations}(a), we vary the KL regularization strength $\beta$ from $0.0$ to $3.0$. Because we initialize the student's merging logits with the MI-guided routing prior at the start of each test batch, the model's parameters are refreshed, preventing long-term drift from accumulating across batches even when $\beta=0.0$. However, when the stream is adapted continuously without frequent routing-prior resets, the teacher regularization becomes crucial to maintain model sanity.

\textbf{Effect of EMA Momentum ($\alpha_{\text{ema}}$):}
We study the impact of the EMA momentum parameter $\alpha_{\text{ema}}$ on model performance in Table~\ref{tab:ablations}(b). Setting $\alpha_{\text{ema}}=0.0$ corresponds to a teacher model that matches the student's weights from the previous step without temporal smoothing. As $\alpha_{\text{ema}}$ increases, the temporal stability of the teacher improves, providing small but consistent accuracy improvements.

\textbf{Effect of Mutual Information Temperature ($\gamma$):}
The scaling parameter $\gamma$ governs the sharpness of the routing prior constructed from the expert mutual informations. Table~\ref{tab:ablations}(c) shows that when $\gamma=1.0$ (weak scaling), the routing prior is near-uniform, failing to correctly route inputs and resulting in poor performance. As $\gamma$ increases to $15.0$ and above, the router clearly isolates the correct expert, establishing near-perfect routing.

\textbf{Effect of Inner Adaptation Steps:}
To evaluate the long-term stability and resilience to the feedback loop trap, we compare AdaMerging and our MI-FRTR-TTMM framework as the number of test-time inner adaptation steps increases from 1 to 10 (Table~\ref{tab:inner_steps}). As the number of adaptation steps increases, unconstrained entropy minimization (\textbf{AdaMerging}) remains stuck in a severely collapsed state, yielding poor performance (around 0.48--0.51 accuracy on Closed Sequential). This empirically confirms Proposition~3.2: unconstrained gradients rapidly drag the model parameters into the OOD expert's overconfident basin. Conversely, our \textbf{MI-FRTR-TTMM} framework is remarkably stable, with its accuracy steadily rising from 0.8781 (at 1 step) to \textbf{0.8848} (at 10 steps) on Closed Sequential, and from 0.6322 to \textbf{0.6359} on Noisy Open-World. The EMA teacher provides a powerful temporal regularizer that ensures optimization remains robust and productive over multiple steps.

\begin{table}[h]
\caption{Comparison of classification accuracies under varying number of test-time inner adaptation steps on the Closed Sequential and Noisy Open-World streams.}
\label{tab:inner_steps}
\centering
\begin{small}
\begin{sc}
\begin{tabular}{ccccc}
\toprule
& \multicolumn{2}{c}{\textbf{Closed Seq}} & \multicolumn{2}{c}{\textbf{Noisy OW}} \\
\textbf{Steps} & \textbf{Ada} & \textbf{Ours} & \textbf{Ada} & \textbf{Ours} \\
\midrule
1  & 0.4805 & 0.8781 & 0.3850 & 0.6322 \\
2  & 0.4672 & 0.8789 & 0.3694 & 0.6322 \\
3  & 0.4820 & 0.8801 & 0.3709 & 0.6331 \\
5  & 0.4984 & 0.8809 & 0.3837 & 0.6341 \\
10 & 0.5152 & \textbf{0.8848} & 0.4028 & \textbf{0.6359} \\
\bottomrule
\end{tabular}
\end{sc}
\end{small}
\vspace{-10pt}
\end{table}

\begin{table*}[t]
\caption{Ablation studies on the core hyperparameters of MI-FRTR-TTMM, evaluated on the Closed Sequential and Noisy Open-World streams. We analyze (a) the teacher regularization strength $\beta$, (b) the EMA teacher momentum $\alpha_{\text{ema}}$, and (c) the Mutual Information scaling parameter $\gamma$.}
\label{tab:ablations}
\centering
\begin{small}
\begin{sc}
\begin{tabular}{ccc|ccc|ccc}
\toprule
\multicolumn{3}{c}{\textbf{(a) Regularization strength $\beta$}} & \multicolumn{3}{c}{\textbf{(b) EMA momentum $\alpha_{\text{ema}}$}} & \multicolumn{3}{c}{\textbf{(c) MI scaling parameter $\gamma$}} \\
$\beta$ & Closed Seq & Noisy OW & $\alpha_{\text{ema}}$ & Closed Seq & Noisy OW & $\gamma$ & Closed Seq & Noisy OW \\
\midrule
0.0 & 0.8805 & 0.6338 & 0.00 & 0.8793 & 0.6325 &  1.0 & 0.7684 & 0.5772 \\
0.5 & 0.8805 & 0.6338 & 0.50 & 0.8797 & 0.6328 &  5.0 & 0.8520 & 0.6228 \\
1.0 & 0.8801 & 0.6334 & 0.80 & 0.8801 & 0.6331 & 10.0 & 0.8711 & 0.6316 \\
1.5 & \textbf{0.8801} & \textbf{0.6331} & 0.90 & \textbf{0.8801} & \textbf{0.6331} & 15.0 & \textbf{0.8801} & \textbf{0.6331} \\
3.0 & 0.8797 & 0.6328 & 0.99 & 0.8801 & 0.6331 & 30.0 & 0.8836 & 0.6331 \\
\bottomrule
\end{tabular}
\end{sc}
\end{small}
\vspace{-10pt}
\end{table*}

\section{Discussion, Limitations, and Conclusion}
In this work, we identified and resolved two fundamental bottlenecks in Test-Time Model Merging: the autograd disconnection bug in prior baselines and the failure of sample-level entropy routing due to OOD overconfidence. By introducing an information-theoretic Mutual Information routing prior and a feedback-resistant EMA teacher with task-switch resetting, our proposed MI-FRTR-TTMM framework establishes a robust, highly performing paradigm for dynamic multi-task inference.

\subsection{Limitations and Broader Impact}
While our framework achieves state-of-the-art results across non-stationary streams, we acknowledge several limitations that provide valuable avenues for future work:

\textbf{Memory Overhead:} FRTR-TTMM maintains an Exponential Moving Average (EMA) teacher model in memory alongside the active student model. While SimpleCNN parameters are extremely lightweight, maintaining duplicate parameters for massive-scale architectures (e.g., billions of parameters LLMs or heavy ViTs) increases memory consumption during test-time adaptation.

\textbf{Sensitivity to Batch Size:} Mutual Information estimation relies on computing marginal prediction entropy $H(Y)$ across a batch of inputs. Under extremely small batch sizes (e.g., $B \leq 8$), the empirical marginal distribution $\bar{p}(y)$ may exhibit higher variance, potentially reducing the accuracy of the routing prior.

\textbf{Routing Inference Cost:} To calculate the Mutual Information routing prior, our method passes the test batch through all $K$ expert models. While this cost is modest for small expert libraries, it scales linearly with the number of experts $K$. Future research will explore sparse, hierarchical, or early-exit routing networks to mitigate this overhead.

\subsection{Conclusion}
Despite these limitations, MI-FRTR-TTMM addresses the fundamental feedback loop trap of test-time adaptation without requiring complex or computationally expensive second-order preconditioning (e.g., Fisher Information). It provides a clean, information-theoretic and temporal consistency framework that is highly effective under severe noise and covariate shifts. Future work will investigate scaling this framework to large multi-modal architectures and complex open-world environments.

\nocite{*}
\bibliography{submission}
\bibliographystyle{icml2026}

\end{document}
